\documentclass[11pt,a4paper]{article}
\usepackage[margin=22mm,headheight=15pt]{geometry}
\usepackage{fontspec}
\setmainfont{lmroman10-regular.otf}[BoldFont=lmroman10-bold.otf,ItalicFont=lmroman10-italic.otf,BoldItalicFont=lmroman10-bolditalic.otf]
\setsansfont{lmsans10-regular.otf}[BoldFont=lmsans10-bold.otf]
\setmonofont{lmmono10-regular.otf}[Scale=MatchLowercase]
\usepackage[ngerman]{babel}
\usepackage{amsmath,xurl}
\usepackage{microtype,booktabs,tabularx,array,enumitem,xcolor,fancyhdr}
\usepackage{hyperref}
\definecolor{accent}{HTML}{176458}
\definecolor{muted}{HTML}{52606B}
\hypersetup{colorlinks=true,urlcolor=accent,linkcolor=accent,pdftitle={Fusionsbericht: So funktioniert der Dataset-Branch},pdfauthor={ETHack}}
\pagestyle{fancy}\fancyhf{}
\fancyhead[L]{\small\sffamily ETHack | Dataset einfach erklärt}
\fancyhead[R]{\small\sffamily Stand 12.09.2026}
\fancyfoot[R]{\small\thepage}
\setlength{\emergencystretch}{2em}
\renewcommand{\arraystretch}{1.15}
\setlength{\parindent}{0pt}\setlength{\parskip}{6pt}
\setlist[itemize]{leftmargin=*,itemsep=3pt,topsep=3pt}
\setlist[enumerate]{leftmargin=*,itemsep=5pt,topsep=3pt}
\newcolumntype{Y}{>{\raggedright\arraybackslash}X}
\newcommand{\pf}[1]{\path{#1}}
\newcommand{\lead}[1]{{\large\sffamily\color{accent}#1}\par}
\newcommand{\page}[2]{\clearpage\section*{#1}\lead{#2}}
\begin{document}
\thispagestyle{empty}
{\sffamily\color{accent}\large ETHack 2026\par}
\vspace{8mm}
{\sffamily\bfseries\Huge Fusionsbericht\par}
\vspace{3mm}
{\sffamily\LARGE So funktioniert der\par Dataset-Branch\par}
\vspace{8mm}
\lead{Daten sammeln. Firmen zuordnen. Werte vergleichen. Grenzen sichtbar machen.}
Das Projekt sammelt Informationen über Firmen im S\&P 500: etwa Emissionen, Arbeitsunfälle, Umweltverstösse und Klimaziele. Es verbindet diese Angaben mit den Firmen und zeigt zu jedem Wert eine Quelle. Für einen kleineren Teil der Firmen berechnet es zusätzlich eine relative Klimabewertung.

\textbf{Das wichtigste Ergebnis:} Viele Daten zu einer Firma zu haben bedeutet noch nicht, ihre gesamte Nachhaltigkeit beurteilen zu können. Die berechnete Kernnote nutzt hier nur drei emissionsbezogene Kennzahlen. Soziales, Governance und Zielangaben stehen daneben.

\vspace{3mm}
\begin{tabularx}{\linewidth}{@{}rY@{}}\toprule
\textbf{503} & Firmen beziehungsweise Indexeinträge in der Firmenansicht\\
\textbf{467} & mit mindestens einem Nachhaltigkeitswert; 36 ohne einen solchen Wert\\
\textbf{143} & mit einem gespeicherten Rangband aus der Klimabewertung\\
\textbf{45} & unterschiedliche Kennzahlen im Datensatz\\
\textbf{11\,798} & Datenzeilen aus 14 verwendeten Quellen; das Register enthält 19 Quellen\\\bottomrule
\end{tabularx}

\vspace{3mm}
\textbf{Welche Fassung ist das?} Diese Erklärung bezieht sich auf den Code und die gespeicherten Ergebnisse im Branch \texttt{dataset}, Ausgangscommit \texttt{aaf65d2}. Sie übernimmt keine Ergebniszahlen ungeprüft aus älteren Berichten. Die Daten wurden für diese Dokumentation gelesen, nicht neu von den Behörden abgerufen oder neu bewertet.

Das Original \texttt{fusionsbericht (1).pdf} wurde unverändert aus \texttt{main}, Commit \texttt{662d939}, in diesen Branch übernommen. Es beschreibt eine weitergehende Modellvariante. Deren Skripte \texttt{f01} bis \texttt{f07} liegen hier nicht vor. Seite 6 zeigt die Unterschiede.

\vfill
{\small\color{muted}Leseroute: Ablauf S. 2 | Bewertung S. 3 | Datensatz S. 4 | Nutzung S. 5 | Abgleich mit dem Original S. 6 | Verbesserungen S. 7--8.}

\page{1. Vom Messwert zur Firmenansicht}{Der Weg besteht aus fünf verständlichen Schritten.}
\begin{enumerate}
\item \textbf{Quellen einlesen.} Die Module unter \pf{pipeline/sources/} laden Behörden- und Organisationsdaten. Team-Dateien kommen aus \pf{data/external/team/}. Das Skript \pf{scripts/01_fetch.py} startet die registrierten Quellen. Es speichert lokale Kopien als Parquet in \pf{data/raw/}; vorhandene Kopien werden ohne \texttt{--force} wiederverwendet.
\item \textbf{Anlagen einer Firma zuordnen.} Eine Behörde kennt oft einen Kraftwerksnamen oder eine Tochterfirma, aber kein Börsenkürzel. \pf{pipeline/resolve.py} bereinigt Namen und sucht zuerst exakte Treffer, dann Einträge in einer Tochtertabelle, danach passende Namensanfänge und schliesslich ähnliche Namen. Der unscharfe Abgleich verwendet standardmässig die Schwelle 92.
\item \textbf{Firmenwerte berechnen.} \pf{pipeline/canonical.py} summiert zugeordnete Emissionen je Firma und Jahr. Die SEC-Kennung CIK verbindet Umsätze mit der Firma. Daraus entstehen Emissionen je Umsatz und mehrjährige Trends. \pf{scripts/02_analyse.py} berechnet die Rangbänder.
\item \textbf{Weitere Themen ergänzen.} \pf{scripts/05b_firmenaggregate.py} verdichtet Stromerzeugung, Giftstoffmeldungen, Arbeitsunfälle und Umweltverstösse auf Firmenebene. Die Regeln der Zuordnung sind nicht überall dieselben wie im Emissionsranking.
\item \textbf{Alles zusammenführen und anzeigen.} \pf{pipeline/consolidate.py} verbindet Ergebnisse, Zusatzquellen und Team-Daten. \pf{scripts/07_dataset.py} schreibt CSV, Parquet und JSON und kopiert die Dateien in die Web-App. Dort lassen sich Firmen, Kennzahlen und Quellen durchsuchen.
\end{enumerate}

\subsection*{Ein Rechenbeispiel}
Eine Anlage meldet 100\,000 Tonnen Emissionen. Firma A besitzt 60\,\% der Anlage. Im GHGRP-Pfad werden ihr 60\,000 Tonnen zugerechnet. Hat sie 2\,000 Mio. USD Umsatz, ergibt das:
\[
\text{Emissionsintensität} = \frac{60\,000\ \text{Tonnen}}{2\,000\ \text{Mio. USD}} = 30\ \text{Tonnen je Mio. USD}.
\]
Das ist ein erfundenes Beispiel zum Verständnis. Die Intensität beantwortet, wie viel Ausstoss mit einer Umsatzmenge verbunden ist. Sie sagt weder, wie klimafreundlich ein Produkt ist, noch ob die gesamte Firma nachhaltig arbeitet.

\textbf{Wichtige Trennung:} Der GHGRP-Code nimmt festgelegte Direktemittenten-Sektoren auf und weist biogenes CO$_2$ separat aus. Lieferanten- und Injektionszeilen dürfen nicht einfach zur direkten Emission addiert werden. Der Dataset-Code nutzt dafür eine feste Liste von Sektor-IDs; deren Vollständigkeit braucht eine eigene Prüfung.

\page{2. Wie die Bewertung funktioniert}{Drei Fragen zur Klimawirkung, viele mögliche Rechenwege.}
\begin{tabularx}{\linewidth}{@{}p{43mm}Y@{}}\toprule
Kennzahl & Einfache Bedeutung\\\midrule
CO$_2$-Intensität & Wie viele Tonnen entstehen je Million USD Umsatz?\\
Trend der Intensität & Wird dieser Wert über mehrere Jahre kleiner oder grösser?\\
Trend der Tonnen & Sinkt der gesamte erfasste Ausstoss oder wächst er?\\\bottomrule
\end{tabularx}
Bei allen drei Kennzahlen gilt in der Bewertung: kleiner ist besser. Die Trends werden aus einer Linie durch die logarithmierten Jahreswerte abgeleitet. Dafür braucht es mindestens drei positive, gültige Jahreswerte. Das verwendete Zeitfenster ist 2018 bis 2023.

\subsection*{Warum Firmen innerhalb einer Branche verglichen werden}
Ein Stromerzeuger und ein Softwareunternehmen haben sehr unterschiedliche Tätigkeiten. Die Bewertung vergleicht deshalb innerhalb von GICS-Sektoren oder deren Unterbranchen. Hat eine Unterbranche weniger als sechs Firmen, fällt der Code auf den Sektor zurück. \textbf{Ein Sektor selbst kann trotzdem weniger als sechs Firmen haben.}

Die drei Kennzahlen werden zuerst auf vergleichbare Skalen gebracht. Danach erhalten sie Gewichte. Schliesslich werden sie zu einer Note zusammengeführt. Der Code probiert sowohl einen gewichteten Durchschnitt als auch ein geometrisches Mittel. Beim geometrischen Mittel ziehen schwache Teilwerte das Ergebnis stärker nach unten; ein vollständiges Ausgleichsverbot ist es nicht.

\subsection*{Warum am Ende ein Band statt eines festen Platzes steht}
Die Hauptauswertung rechnet 1\,500 zufällige Durchläufe. Sie variiert vier Skalierungsverfahren, drei Gewichtungen, zwei Arten der Zusammenfassung, zwei Vergleichsebenen und drei Einstellungen für Extremwerte. Manchmal wird eine der drei Kennzahlen weggelassen. Dazu kommen kleine zufällige Änderungen der Eingangswerte, stärker bei unsicherer Zuordnung.

Die diskreten Optionen ergeben nominell 576 Kombinationen. Die 1\,500 Durchläufe sind eine Zufallsstichprobe mit Wiederholungen, kein vollständiges Experiment mit gleich vielen Durchläufen pro Kombination. Der Zufallsstartwert ist \texttt{20260912}.

\textbf{Beispiel: Band 30 bis 80, Mitte 60.} Unter den gewählten Annahmen liegt die Firma oft im Mittelfeld, die Einordnung schwankt aber stark. Höhere Perzentile bedeuten bessere relative Ergebnisse. Das Band reicht vom 10. bis zum 90. Perzentil der Durchläufe. Es ist kein Beweis, dass der wahre Rang mit 80\,\% Wahrscheinlichkeit darin liegt.

Bei den 143 gespeicherten Firmen ist das Band im Median 34,5 Perzentilpunkte breit. Zwischen Branchen sind gleiche Perzentile keine gleichen absoluten Umweltleistungen. Die Zusatzdaten zu Unfällen, Zielen oder Verstössen verändern diese Kernnote derzeit nicht.

\page{3. Was im Datensatz steckt}{Eine Zeile enthält einen Wert mit Firma, Jahr und Herkunft.}
\pf{data/out/dataset_long.csv} und die Parquet-Datei enthalten dieselbe Art von Tabelle: eine Zeile je Börsenkürzel, Jahr und Kennzahl. So lassen sich mehrere Jahre und verschiedene Quellen zusammen speichern, ohne alles in eine einzige grosse Firmenzeile zu pressen.

\begin{tabularx}{\linewidth}{@{}p{47mm}Y@{}}\toprule
Feld & Bedeutung\\\midrule
\texttt{ticker}, \texttt{company} & Börsenkürzel und Firmenname\\
\texttt{year} & derzeit ein einzelnes zugewiesenes Jahr; bei manchen Werten nur eine vereinfachte Periodenangabe\\
\texttt{metric}, \texttt{value}, \texttt{unit} & Kennzahl, Zahlenwert und Einheit\\
\texttt{axis}, \texttt{direction} & Themenbereich; ob kleinere, grössere oder keine Werte bevorzugt werden\\
\texttt{source\_id}, \texttt{source\_url} & Quelle und allgemeine Abrufadresse\\
\texttt{source\_license} & im Register hinterlegte Nutzungsangabe\\
\texttt{retrieved\_at} & im aktuellen Export das Erzeugungsdatum, nicht verlässlich das ursprüngliche Abrufdatum\\
\texttt{partial\_year} & Kennzeichnung des laufenden CAMPD-Jahres\\\bottomrule
\end{tabularx}

\subsection*{Die Themenbereiche in einfacher Sprache}
\begin{itemize}
\item \textbf{A: Umweltwerte.} Emissionen, Stromintensität und Giftstofffreisetzungen. Nur die drei Kennzahlen auf Seite 3 gehen in die Kernnote ein.
\item \textbf{B: Ziele und Warnsignale.} Zum Beispiel SBTi-Ziele und der Hinweis, dass die Intensität sinkt, während die Tonnen steigen. Das allein belegt keine Täuschung.
\item \textbf{S: Soziales.} Arbeitsunfälle, Lohnverfahren und ergänzende WBA-Bewertungen.
\item \textbf{G: Umwelt-Compliance.} ECHO-Angaben zu Strafen und Verstössen. Das ist ein Ausschnitt; Vorstandskontrolle oder Vergütung werden damit nicht abgedeckt.
\item \textbf{Metadaten, Ergebnis, Vergleich.} Umsätze und Anlagenzahlen, berechnete Rangbänder sowie ein kommerzieller Vergleichswert. Der Vergleichswert fliesst nicht in die eigene Note ein.
\end{itemize}

\textbf{Fehlend ist nicht null.} Kein Treffer kann heissen, dass eine Firma nicht erfasst wurde oder die Zuordnung scheiterte. Ein registrierter Quellenname beweist ausserdem nicht, dass diese Quelle im Export tatsächlich Werte liefert: 19 Registereinträgen stehen 14 verwendete Quellen gegenüber.

\page{4. Wie man die Ergebnisse benutzt}{Immer Firma, Einheit, Zeitraum und Quelle gemeinsam lesen.}
Die Web-App liest die vorbereiteten Dateien unter \pf{web/public/data/}. Sie fragt beim Anklicken einer Firma keine Behörde live ab. \texttt{companies.json} enthält je Kennzahl bevorzugt den jüngsten vollständigen Wert. Ein laufendes CAMPD-Jahr wird nur genommen, wenn kein vollständiger Wert verfügbar ist. Dadurch können die Werte derselben Firma aus unterschiedlichen Jahren stammen.

\begin{tabularx}{\linewidth}{@{}p{42mm}Y@{}}\toprule
Ansicht & Wofür sie hilfreich ist\\\midrule
Firmenliste und Firmenprofil & Firma suchen, Branche filtern, einzelne Werte mit Quelle ansehen\\
Kennzahlen und Quellen & Bedeutung, Einheit und Datenabdeckung nachschlagen\\
Vergleich & Firmen anhand ausgewählter Kennzahlen nebeneinander betrachten\\
Downloads & Den vorbereiteten Datensatz und seine Register weiterverwenden\\\bottomrule
\end{tabularx}

\textbf{Sinnvolle Lesereihenfolge:} Zuerst die verfügbare Kennzahl und das Jahr prüfen. Dann Firmen mit ähnlicher Tätigkeit vergleichen. Danach die Datenlücken ansehen. Ein Rangband erst zum Schluss als Zusammenfassung der drei Klimaindikatoren lesen.

\subsection*{Gespeicherte Ergebnisse sofort ansehen}
Vom Repository-Verzeichnis aus, mit installiertem Node.js/npm:
\begin{verbatim}
cd web
npm ci
npm run dev
\end{verbatim}
Danach \texttt{http://localhost:3000} öffnen. Dafür reichen die bereits gespeicherten Web-Daten.

\subsection*{Daten neu aufbereiten}
Mit eingerichteter Python-Umgebung, benötigten Paketen und verfügbaren Rohdateien laufen vom Repository-Verzeichnis aus:
\begin{verbatim}
python scripts/02_analyse.py
python scripts/05b_firmenaggregate.py
python scripts/07_dataset.py
\end{verbatim}
Der erste Schritt berechnet das Ranking, der zweite die Zusatzaggregate, der dritte die Exporte und Web-Daten. \texttt{07\_dataset.py} allein berechnet kein neues Ranking. Für frische Rohdaten geht \texttt{01\_fetch.py} voraus; einzelne Quellen benötigen Schlüssel oder lokale Team-Dateien. Es meldet Quellenfehler und läuft dann weiter. Ein erfolgreicher Programmabschluss garantiert deshalb keinen vollständigen Abruf.

Die Python-Abhängigkeiten für den Grundpfad stehen im README; einzelne Zusatzquellen brauchen weitere Pakete. Rohdaten werden nicht mit Git versioniert. Eine vollständige Reproduktion aus einem frischen Checkout ist deshalb noch kein einzelner, abgesicherter Befehl.

\page{5. Was aus dem Original übernommen wurde}{Das Original ist verfügbar; seine Modellvariante ist noch nicht integriert.}
Der Fusionsbericht aus \texttt{main} beschreibt eine eigenständige Weiterentwicklung. Die Tabelle trennt seine Aussagen vom tatsächlich vorliegenden Dataset-Stand. Die Originalzahlen sind hier als Berichtsaussagen wiedergegeben, nicht neu reproduziert.

\small
\begin{tabularx}{\linewidth}{@{}p{31mm}YY@{}}\toprule
Thema & Original auf \texttt{main} & Vorliegender Dataset-Branch\\\midrule
Skripte & \texttt{f01\_sources.py} bis \texttt{f07\_supplement.py} & \texttt{scripts/01} bis \texttt{08}, Module unter \texttt{pipeline/}; die \texttt{f}-Skripte fehlen\\
Kernnote & Sechs Klimaindikatoren, später ein siebter zu Giftstoffen & Drei Kennzahlen zu Emissionsintensität und Trends\\
Vergleichsgruppen & Aus Emissionsprofilen gebildete Gruppen plus GICS & GICS-Sektor oder Unterbranche\\
Unsicherheit & 768 Konfigurationen, je 100 Ziehungen & 1\,500 zufällige Durchläufe aus nominell 576 Kombinationen plus Rauschen\\
Zuordnung & Vier Schwellen, zwei Zurechnungsregeln & Hauptpfad: Schwelle 92 und Beteiligungsanteile; Zusatzquellen mit anderen Regeln\\
Rangbänder & 116 Firmen, mediane Breite 52,9 Punkte & 143 Firmen, mediane Breite 34,5 Punkte\\
Soziales und Governance & Unfallachse, Governance als offen beschrieben & Zusätzliche Sozialdaten und ECHO-Compliance-Werte im Dataset, ausserhalb der Kernnote\\
Portfolio & Beschreibt eine Portfolio-Optimierung & Dieser Rechenpfad ist hier nicht vorhanden\\\bottomrule
\end{tabularx}
\normalsize

\subsection*{Was man daraus schliessen darf}
Die gemeinsame Idee bleibt sinnvoll: Daten nachvollziehbar zuordnen, ähnliche Firmen vergleichen und Unsicherheit sichtbar machen. Der Dataset-Branch ergänzt dazu einen breiten Datenbestand und eine Oberfläche.

\textbf{Was noch fehlt, ist die Verbindung der weitergehenden Methodik mit diesem Code.} Das Kopieren einer PDF übernimmt weder deren Programme noch die zugehörigen Daten. Auch ein schmaleres Band beweist keine bessere Methode, wenn andere Firmen, Kennzahlen oder Annahmen verwendet werden.

Im Original stehen Aussagen wie vollständig fehlende Governance oder fehlende physische Nenner. Für diesen Branch muss man sie genauer fassen: Es gibt ECHO-Compliance-Werte für 136 Firmen und Stromintensitäten aus eGRID für 45 Firmen. Diese Angaben ergeben noch keine umfassende Governance-Bewertung und sind nicht in die Kernnote eingebaut.

\page{6. Was zuerst verbessert werden muss}{Zuerst die Verlässlichkeit einzelner Werte, dann ein grösseres Ranking.}
\subsection*{Priorität 1: Konzernzuordnung mit Belegen prüfen}
\textbf{Problem:} Die Tochtertabelle und Namensanfänge entscheiden über grosse Emissionsmengen. Zusatzquellen nutzen unterschiedliche Regeln: GHGRP gewichtet Eigentumsanteile; CAMPD rechnet jedem gefundenen Eigentümer die volle Anlagenmenge zu. Bei mehreren Index-Eigentümern sind so Mehrfachzurechnungen möglich. eGRID ordnet über Betreiber beziehungsweise Versorger zu, ECHO über den zuerst genannten Eigentümer.

\textbf{Verbesserung:} Eine gemeinsame Anlage-Tochter-Mutter-Tabelle mit Gültigkeitsjahr, Rolle, Anteil und Beleg führen. Zunächst die grössten Emissionsbeiträge und unklaren Treffer manuell prüfen. Danach Summen und Eigentumsanteile automatisch kontrollieren. Eine Namensähnlichkeit von 0,95 ist keine gemessene 95-prozentige Trefferwahrscheinlichkeit.

\subsection*{Priorität 2: Einheiten, Zeiträume und Herkunft richtig speichern}
\textbf{Problem:} Die Trendfunktion liefert Dezimalraten, etwa \texttt{0.05} für 5\,\%. Im Export steht dafür die Einheit \texttt{\%/Jahr}, ohne Umrechnung. Die Web-Anzeige rechnet korrekt mit Faktor 100 um; wer nur die CSV liest, kann den Wert falsch verstehen. Mehrjahreswerte werden auf ein Jahr verkürzt: Lohnverfahren 2022--2024 etwa auf 2024. ECHO bekommt pauschal 2025. Beim Export wird das aktuelle Datum als \texttt{retrieved\_at} gesetzt, auch bei alten lokalen Daten.

\textbf{Verbesserung:} Eine eindeutige Prozentkonvention wählen und im Export prüfen. Start- und Enddatum, Berichtsjahr, echtes Abrufdatum und Exportdatum getrennt speichern. Originaldatei, Zeilenkennung und Berechnungsweg bis zum Ergebnis erhalten. Allgemeine Quellenlinks reichen für eine Prüfung einzelner Werte noch nicht aus.

\subsection*{Priorität 3: Fehlende Daten dürfen keine sichere Note erzeugen}
\textbf{Problem:} Die Rangliste hat 143 Einträge. Im Export gibt es aber nur für 130 Firmen einen absoluten Trend und für 127 einen Intensitätstrend. Drei beobachtete Jahre garantieren keine drei gültigen positiven Werte. Die Zusammenfassung kann über die verbleibenden Kennzahlen rechnen; bei konstanten oder leeren Gruppen kann die Skalierung einen neutralen Wert erzeugen. Ein Vergleich mit fehlenden Trends kann ausserdem ein Warnsignal als \texttt{false} ausgeben, obwohl es unbekannt ist.

\textbf{Verbesserung:} Mindestdaten je Kennzahl prüfen. Vollständige Bewertung, Teilbewertung und nicht bewertbar ausdrücklich trennen. Fehlend, null und falsch als verschiedene Zustände führen. Kleine Sektorgruppen sperren oder begründet zusammenlegen. Die 467 Firmen mit irgendeinem Nachhaltigkeitswert dürfen nicht als 467 vollständig bewertete Firmen erscheinen.

\page{7. Was danach den grössten Nutzen bringt}{Methodik erweitern und jeden Bericht an einen überprüfbaren Datenstand binden.}
\subsection*{Priorität 4: Abdeckung und Vergleichbarkeit verbessern}
Die Kernnote betrachtet erfasste US-Anlagen. Zugekaufter Strom (Scope 2), vor- und nachgelagerte Aktivitäten (Scope 3) sowie Auslandsanlagen sind damit nicht vollständig erfasst. Umwelt-, Sozial- und Governance-Themen bleiben sehr unterschiedlich abgedeckt. Fehlende freie Werte sind hier eine Lücke im Bestand; daraus folgt nicht, dass solche Daten grundsätzlich nirgends erhältlich sind.

\textbf{Nächster Schritt:} Für eine klar abgegrenzte Branche zusätzliche Quellen und denselben Berichtszeitraum zusammenführen. Bei Stromerzeugern die vorhandenen MWh-Daten verwenden, sobald Emissionen, Erzeugung und Eigentumsgrenzen zusammenpassen. US-Anlagenemissionen durch weltweiten Konzernumsatz zu teilen kann die Vergleichbarkeit verzerren.

\subsection*{Priorität 5: Die Methodik aus dem Original reproduzierbar integrieren}
Die beiden Trendkennzahlen hängen zusammen: In \pf{data/out/kennzahlen.json} beträgt ihre gespeicherte Rangkorrelation rund 0,59. Sie liefern also teilweise ähnliche Information. Die zufälligen Störungen beruhen zudem auf gesetzten Parametern und unkalibrierten Zuordnungswerten.

\textbf{Nächster Schritt:} Die im Original genannten Skripte und Eingabestände beschaffen, getrennt reproduzieren und erst dann übernehmen. Mit identischem Firmenkreis prüfen, ob neue Kennzahlen eigenständige Information liefern und neue Vergleichsgruppen sinnvoll sind. Methodenunsicherheit und Datenfehler getrennt berichten; Rauschannahmen mit geprüften Beispielen kalibrieren.

\subsection*{Priorität 6: Aktualisierung und Dokumentation absichern}
\pf{01_fetch.py} setzt den Lauf trotz Quellenfehlern fort. \pf{consolidate.py} überspringt einige fehlende Dateien und behält bei doppeltem Schlüssel still den ersten Wert. Ein alter Rang kann deshalb zusammen mit neueren Zusatzdaten veröffentlicht werden. README, Quellenregister und ältere Berichte enthalten zudem abweichende Abdeckungszahlen.

\textbf{Nächster Schritt:} Pro Lauf ein Manifest mit Quellenversionen, Zeiträumen, Dateiprüfsummen und Fehlerstatus schreiben. Bei fehlenden Pflichtquellen oder widersprüchlichen Doppelwerten abbrechen. Abdeckung, Einheiten und Rangzulassung vor dem Export prüfen. Berichtskennzahlen automatisch aus genau diesem Lauf erzeugen und Abhängigkeiten festhalten.

\subsection*{Prüfgrundlage dieser Fassung}
{\small Gelesen wurden das Original-PDF aus \texttt{main}, die auf den vorherigen Seiten genannten Python-Module, die Web-Datenlogik und die gespeicherten CSV-/JSON-Dateien. Die Übersichtszahlen wurden direkt nachgezählt; der Langdatensatz enthält keine doppelten Schlüssel aus Firma, Jahr und Kennzahl. Das ist eine Bestandsprüfung, keine unabhängige Bestätigung aller Ursprungswerte.}

{\small PDF neu bauen, vom Repository-Verzeichnis aus mit Tectonic 0.15.0:\par
\texttt{tectonic --untrusted --outdir report report/fusionsbericht.tex}\par
Anschliessend Compilerwarnungen prüfen und Seiten mit \texttt{pdfinfo} und \texttt{pdftoppm} kontrollieren. Die Quellen dieser Erklärung sind im Repository prüfbar; externe Verfügbarkeits- und Rechtsaussagen wurden nicht neu recherchiert.}
\end{document}
